\documentclass{ctexart}
\usepackage{amsmath, amssymb, amsthm}
\usepackage{geometry}
\usepackage{enumitem}
\usepackage{hyperref}
\usepackage{boxedminipage}

\newtheorem{remark}{评注}

\begin{document}

整体思路是对的，尤其是这句话：

\begin{quote}
Geometric theory 用来构造理论的分类对象；sheaf semantics 用来在这些分类对象上解释和推理。
\end{quote}

但有几处需要严格修正。最重要的是：

\[
\boxed{
\text{命题 geometric theory 通常分类 locale；一般 geometric theory 分类 topos。}
}
\]

不能把这两种情况完全混同。

\section*{1. 命题理论对应 classifying locale：基本正确}

如果 \(\mathbb T\) 是命题 geometric theory，那么它没有 sort、变量和带参数的关系符号，只有命题符号以及

\[
\top,\quad\bot,\quad\wedge,\quad\bigvee.
\]

其公式模可证明等价形成 frame：

\[
\mathcal O(\mathbb T)
=
\{\text{命题 geometric formulas}\}/\mathord{\dashv\vdash_{\mathbb T}}.
\]

因此得到一个 locale：

\[
L(\mathbb T).
\]

在这种情形下：

\[
\boxed{
\text{\(\mathbb T\) 的 classifying object 是一个 locale。}
}
\]

并且它的点对应 \(\mathbb T\) 的 set-based models：

\[
\operatorname{Pt}(L(\mathbb T))
\cong
\operatorname{Mod}_{\mathbf{Set}}(\mathbb T).
\]

更具体地，一个 set-based 模型就是一个真值赋值

\[
v:\{\text{命题符号}\}\to\Omega
\]

并满足所有公理；这等价于 Lindenbaum frame 上的 frame homomorphism

\[
\mathcal O(\mathbb T)\to\Omega,
\]

也就是完全素滤子。

所以你第 1 点的核心说法是正确的。

不过，``任意一个命题 geometric theory 都对应一个 classifying locale''最好补充：

\begin{quote}
这里的 classifying locale 是命题理论的 classifying topos 的局部对象；一般 geometric theory 的分类对象不一定是 locale，而是 classifying topos。
\end{quote}

\section*{2. 一般 geometric theory 的分类对象是 classifying topos}

如果 \(\mathbb T\) 含有 sort 和变量，例如环的理论，那么其模型不是简单的真假赋值。

一个 \(\mathbb T\)-模型 \(M\) 在 \(\mathbf{Set}\) 中包括：

\begin{itemize}
\item 每个 sort 的集合；
\item 每个函数符号的函数；
\item 每个关系符号的子集；
\item 满足所有 geometric sequents。
\end{itemize}

这种理论一般由一个 classifying topos

\[
\mathbf{Set}[\mathbb T]
\]

分类，而不是由一个 locale 分类。

它满足：

\[
\operatorname{Mod}_{\mathbb T}(\mathcal E)
\cong
\operatorname{Geom}(\mathcal E,\mathbf{Set}[\mathbb T])
\]

的自然对应，其中 \(\mathcal E\) 是适当的 topos。

因此应区分：

\[
\boxed{
\begin{aligned}
\text{命题 geometric theory}
&\longrightarrow \text{classifying locale},\\
\text{一般 geometric theory}
&\longrightarrow \text{classifying topos}.
\end{aligned}
}
\]

在某些特殊情况下，一般理论的 classifying topos 也是 localic 的，这时才能进一步由某个 locale 表示。

\section*{3. ``定义 sheaf model''：基本正确，但需要精确化}

在一个 locale \(X\) 上，严格说通常考虑其 sheaf topos：

\[
\operatorname{Sh}(X).
\]

一个 \(\mathbb T\)-模型是 \(\operatorname{Sh}(X)\) 中的一个模型：

\begin{itemize}
\item 每个 sort \(A\) 解释为一个对象
  \[
  A^M\in\operatorname{Sh}(X);
  \]
\item 函数符号解释为态射；
\item 关系符号解释为子对象；
\item 所有公理在内部逻辑中成立。
\end{itemize}

因此，将 sort 说成``sheaf''，在普通空间 \(X\) 上是直观且正确的；在一般 locale 上，更严格地说是 \(\operatorname{Sh}(X)\) 中的对象。

``关系符号解释为子 sheaf''也基本正确，但更一般的说法是：

\[
R^M\hookrightarrow
A_1^M\times\cdots\times A_n^M
\]

是一个子对象。

它未必需要单独称作某种传统意义上的``子 sheaf''。

\section*{4. Generic model：方向需要修正}

你所说的 generic model 基本正确，但最好用几何态射明确写出。

设 \(\mathbb T\) 的 classifying topos 是

\[
\mathbf{Set}[\mathbb T],
\]

其中有一个 generic model

\[
M_{\mathrm{gen}}.
\]

对任意 topos \(\mathcal E\)，每个 \(\mathbb T\)-模型 \(M\) 对应一个几何态射

\[
f:\mathcal E\to\mathbf{Set}[\mathbb T],
\]

并且

\[
M\cong f^\ast(M_{\mathrm{gen}}).
\]

所以准确表述是：

\[
\boxed{
\text{任意 \(\mathcal E\) 中的 \(\mathbb T\)-模型，
都是 generic model 沿某个几何态射的逆像。}
}
\]

如果 \(\mathbb T\) 是命题理论，并且 classifying topos 是

\[
\operatorname{Sh}(L(\mathbb T)),
\]

那么对另一个 locale \(Y\)，一个 \(\mathbb T\)-模型对应于 locale morphism

\[
Y\to L(\mathbb T).
\]

其模型是 generic model 沿对应几何态射拉回得到的。

因此你的第 3 点应改成：

\begin{quote}
对命题 geometric theory，generic model 位于 \(\operatorname{Sh}(L(\mathbb T))\) 中；对一般 geometric theory，它位于 classifying topos \(\mathbf{Set}[\mathbb T]\) 中。
\end{quote}

另外，``沿着某个 locale morphism 拉回''只适用于分类对象和语义对象都处于 localic 情形。一般情况下应说``沿几何态射取逆像''。

\section*{5. 关于 \(\operatorname{Spec}(A)\)：方向基本正确}

给定交换环 \(A\)，可以构造一个命题 geometric theory，其命题符号通常表示基本开集：

\[
D(a),\qquad a\in A.
\]

相应的 frame 由这些生成元以及描述它们关系的 geometric axioms 生成，例如：

\[
D(0)=\bot,
\]

\[
D(1)=\top,
\]

\[
D(ab)=D(a)\wedge D(b),
\]

以及

\[
D(a)\le D(a+b)\vee D(b).
\]

这里最后一个关系反映：

\[
V(a+b)\supseteq V(a)\cap V(b)
\]

或等价的谱空间基本开集关系。

由此得到 Zariski locale：

\[
\operatorname{Spec}(A).
\]

其点对应于 \(A\) 的素理想：

\[
\operatorname{Pt}(\operatorname{Spec}(A))
\cong
\operatorname{Spec}_{\mathrm{set}}(A).
\]

不过，``素滤子理论''这个说法要看你具体采用的 presentation。更常见的说法是：

\begin{itemize}
\item 素理想的 geometric theory；
\item 或基本开集 \(D(a)\) 的命题 geometric theory；
\item 或 \(A\) 的素滤子理论。
\end{itemize}

这些描述可以等价地呈现同一个 Zariski locale，但需要明确具体生成元和公理。

\section*{6. 关于 \(A^\sim\) 是场：需要避免过强表述}

你写道：

\[
\operatorname{Sh}(\operatorname{Spec}(A))
\models
\text{``\(A^\sim\) 是场''}.
\]

这句话在某些构造中是正确的，但需要说明 \(A^\sim\) 到底是什么。

通常在 Zariski locale 上，考虑由 \(A\) 得到的结构 sheaf 或 generic localization。内部可以构造一个 generic prime ideal：

\[
\mathfrak p_{\mathrm{gen}}\subseteq A,
\]

并考虑商环：

\[
A/\mathfrak p_{\mathrm{gen}}.
\]

在适当的内部语义下，这个商环是一个域；或者其进一步局部化得到内部域。这种结论表达了：

\begin{quote}
在 generic prime point 上，环 \(A\) 看起来像一个域。
\end{quote}

但要区分几个对象：

\begin{enumerate}
\item 结构 sheaf \(\mathcal O_{\operatorname{Spec}(A)}\)；
\item constant sheaf associated to \(A\)；
\item generic prime ideal 的商；
\item stalk
   \[
   \mathcal O_{\operatorname{Spec}(A),\mathfrak p};
   \]
\item residue field
   \[
   \kappa(\mathfrak p).
   \]
\end{enumerate}

结构 sheaf本身一般不是一个域，因为其全局截面是 \(A\)，而 \(A\) 不必是域。正确结论通常是关于 generic prime quotient、generic stalk 或 residue-field 型对象，而不是笼统地说``结构 sheaf 是场''。

\section*{7. ``所有 sheaf 都建立在这些 locale 上''过于绝对}

这句话：

\begin{quote}
我们后续所有的 sheaf 和 sheaf semantics 都是建立在这些由 geometric theory 呈现出来的 locale 上的。
\end{quote}

如果是在某篇特定论文的局部框架内，可能是对的；但作为一般数学陈述，它太强。

更准确的是：

\begin{quote}
在命题 geometric theory 的情形，理论可以由 classifying locale 呈现，并可在该 locale 的 sheaf topos 中使用 generic model 和 sheaf semantics。对于一般 geometric theory，分类对象通常是 classifying topos，而不一定是 locale。
\end{quote}

此外，sheaf semantics 不一定必须先从 geometric theory 构造 locale。给定任意 locale \(X\)，都可以直接研究：

\[
\operatorname{Sh}(X)
\]

中的对象及其内部逻辑。Geometric theory 提供的是：

\begin{itemize}
\item 要研究的结构；
\item 公理；
\item generic model；
\item 分类性质；
\item 与其他模型之间的统一对应。
\end{itemize}

它不是 sheaf semantics 存在的前提。

\section*{8. 更准确的工作流}

你的工作流可以改成下面这样：

\subsection*{命题理论情形}

\begin{enumerate}
\item 为目标对象写出命题 geometric theory \(\mathbb T\)；
\item 构造 Lindenbaum frame：
   \[
   \mathcal O(\mathbb T);
   \]
\item 将其看作 classifying locale：
   \[
   L(\mathbb T);
   \]
\item 在
   \[
   \operatorname{Sh}(L(\mathbb T))
   \]
   中得到 generic model；
\item 使用 Kripke–Joyal/sheaf semantics 进行内部推理；
\item 由分类性质把结果传递到任意相应模型。
\end{enumerate}

\subsection*{一般结构理论情形}

\begin{enumerate}
\item 写出带 sorts、函数和关系的 geometric theory \(\mathbb T\)；
\item 构造 classifying topos：
   \[
   \mathbf{Set}[\mathbb T];
   \]
\item 在其中取 generic model；
\item 对任意 topos \(\mathcal E\) 中的模型，通过几何态射拉回；
\item 在 \(\mathcal E\) 的内部逻辑中使用该模型。
\end{enumerate}

\section*{最终评价}

你的四点可以概括为：

\begin{table}[h]
\centering
\begin{tabular}{|p{5cm}|p{4cm}|p{6cm}|}
\hline
\textbf{原说法} & \textbf{评价} & \textbf{需要的修正} \\
\hline
geometric theory 呈现 locale & 对命题理论正确 & 一般理论呈现 classifying topos \\
\hline
geometric theory 定义 sheaf model & 正确 & sheaf 是 \(\operatorname{Sh}(X)\) 中对象，关系是子对象 \\
\hline
存在 generic model & 正确 & 一般沿几何态射拉回，不一定是 locale morphism \\
\hline
用于 \(\operatorname{Spec}(A)\) 和构造性代数 & 正确 & 需具体说明理论、公理及所称的``generic field''对象 \\
\hline
\end{tabular}
\end{table}

因此，最准确的一句话是：

\begin{quote}
\textbf{Geometric theory 是描述数学结构及其局部稳定公理的语法工具；它的分类对象在命题情形下是 locale、一般情形下是 classifying topos。Sheaf semantics 则是在这些分类对象或任意其他 topos 中解释 geometric theory、构造 generic model，并进行局部的构造性推理。}
\end{quote}

\end{document}